\section{Birth templates and scanner scales of primaries}
\label{sec:area_law_primary_birth_templates}

Retain the primary birth regions $R_j$, fragments $F_{j,z}$, and index sets
$J_j$ (Definition~\ref{def:area_law_primary_regions}) at layer $k$ with
$r=2^k$, fine side $t_k$, and belt pitch $s_k$
(Definition~\ref{def:area_law_fine_belt_scales}). A primary is present iff
$J_j\ne\varnothing$, in which case \cite[Section~11]{OpenAI2026AreaLaw} assigns
its scanner scale and birth template as follows.

% Source: 10-geometry.tex, lines 709--729 (geometry:primary-scale and
% geometry:primary-limits) and lines 756--774 (order of the constants).

\begin{definition}[Primary scanner scale and birth template]
    \label{def:area_law_primary_birth_template}
    \lean{TNLean.PEPS.AreaLaw.Geometry.primaryScannerScale,
        TNLean.PEPS.AreaLaw.Geometry.primaryBirthTemplate}
    \leanok
    \uses{def:area_law_primary_regions,def:area_law_fine_belt_scales,
        def:area_law_template}
    For a constant $C_P$, define the primary scanner scale at layer $k$ by
    \begin{align}
        n_k & =\bigl\lceil C_P r^{1+2\delta_0}\bigr\rceil.
        \label{eq:area_law_primary_scanner_scale}
    \end{align}
    For $J_j\ne\varnothing$ and $n\in\mathbb N_{\ge1}$ with
    $C_{\rm tpl}|J_j|(r+1)\le n$ and $2^p\le n$, the primary birth template is
    the template $T$ (Definition~\ref{def:area_law_template}) with parameters
    $(n,s_0)=(n,r)$, pieces $\{F_{j,z}\}_{z\in J_j}$, and points
    $\bigcup_{z\in J_j}(F_{j,z}\cap\mathbb Z^2)$.
\end{definition}

\begin{lemma}[Points and pieces of a primary birth template]
    \label{thm:area_law_primary_birth_template_points}
    \lean{TNLean.PEPS.AreaLaw.Geometry.primaryBirthTemplate_points,
        TNLean.PEPS.AreaLaw.Geometry.primaryBirthTemplate_pieceCount,
        TNLean.PEPS.AreaLaw.Geometry.card_primaryFragmentIndices_le_rpow,
        TNLean.PEPS.AreaLaw.Geometry.card_primaryBirthTemplate_pieces_le}
    \leanok
    The primary birth template $T$ of Definition~\ref{def:area_law_primary_birth_template}
    has $N_0=|J_j|$ pieces with $\operatorname{diam}_\infty(F_{j,z})\le r$,
    $\operatorname{diam}_\infty(T)\le 2^p$, and (at $p=p_k$)
    \begin{align}
        T   & =R_j\cap\mathbb Z^2, &
        N_0 & \le 9r^{2\delta_0}.
        \label{eq:area_law_primary_birth_points_count}
    \end{align}
\end{lemma}
\begin{proof}\leanok
    \uses{def:area_law_primary_birth_template,
        thm:area_law_primary_fragment_union,
        thm:area_law_primary_fragment_rectangles,
        thm:area_law_primary_diameter_separation,
        thm:area_law_primary_fragment_count,
        thm:area_law_fine_belt_scales}
    Each $F_{j,z}$ ($z\in J_j$) has endpoints of the form $o_i+m$ ($m\in\mathbb Z$)
    with strict inequality, so $F_{j,z}\cap\mathbb Z^2\ne\varnothing$; since
    $R_j=\bigcup_{z\in J_j}F_{j,z}$ (Lemma~\ref{thm:area_law_primary_fragment_union}),
    $T=R_j\cap\mathbb Z^2$, and Lemmas~\ref{thm:area_law_primary_fragment_rectangles}
    and~\ref{thm:area_law_primary_diameter_separation} give the diameter bounds.
    With $m=p_k-k\ge0$, $p_k\le(1+\delta_0)k$
    (Lemma~\ref{thm:area_law_fine_belt_scales}) gives $1\le 2^m\le r^{\delta_0}$,
    so Lemma~\ref{thm:area_law_primary_fragment_count} yields
    \begin{align}
        N_0 & \le(2^m+2)^2\le9\cdot4^m\le9r^{2\delta_0}.
        \notag
    \end{align}
\end{proof}

\begin{lemma}[Choice of the primary scale constant]
    \label{thm:area_law_primary_scanner_constant}
    \lean{TNLean.PEPS.AreaLaw.Geometry.primaryBirthTemplate_scale,
        TNLean.PEPS.AreaLaw.Geometry.two_pow_pitchScaleIndex_le_primaryScannerScale,
        TNLean.PEPS.AreaLaw.Geometry.exists_primaryScannerConstant}
    \leanok
    For any $C_{\rm tpl}$, setting $C_P=\max\{18C_{\rm tpl},1\}$ in~(\ref{eq:area_law_primary_scanner_scale})
    ensures that every primary at belt pitch $p_k$ satisfies
    \begin{align}
        n_k                          & \ge C_{\rm tpl}N_0(r+1), &
        \operatorname{diam}_\infty T & \le s_k\le n_k,
        \label{eq:area_law_primary_scanner_constant}
    \end{align}
    so $T$ is a template of Definition~\ref{def:area_law_template} with
    parameters $(n_k,r)$.
\end{lemma}
\begin{proof}\leanok
    \uses{thm:area_law_primary_birth_template_points,
        def:area_law_primary_birth_template}
    For $C_{\rm tpl}>0$ and $r\ge1$,~(\ref{eq:area_law_primary_birth_points_count})
    gives
    \begin{align}
        C_{\rm tpl}N_0(r+1)
         & \le C_{\rm tpl}\cdot9r^{2\delta_0}\cdot2r
        =18C_{\rm tpl}r^{1+2\delta_0}\le n_k,
        \notag
    \end{align}
    and $C_P\ge1$ gives $s_k=2^{p_k}\le r^{1+\delta_0}\le r^{1+2\delta_0}\le n_k$.
\end{proof}

\begin{lemma}[Cut clearance of a primary birth template]
    \label{thm:area_law_primary_birth_template_clearance}
    \lean{TNLean.PEPS.AreaLaw.Geometry.primaryBirthTemplate_isSeparated}
    \leanok
    If $4D_0<C_0$ (for example $C_0>100(D_0+1)$ as in
    \cite[Section~11]{OpenAI2026AreaLaw}), then every primary birth template $T$
    satisfies
    \begin{align}
        \operatorname{dist}_\infty(T,Z) & >4D_0s_0.
        \label{eq:area_law_primary_birth_clearance}
    \end{align}
\end{lemma}
\begin{proof}\leanok
    \uses{thm:area_law_primary_birth_template_points,
        thm:area_law_dyadic_layer_distance}
    Since $T\subseteq R_j\subseteq\overline{D_k}$
    (Lemma~\ref{thm:area_law_primary_birth_template_points}),
    Lemma~\ref{thm:area_law_dyadic_layer_distance} gives
    $\operatorname{dist}_\infty(T,Z)\ge C_0r>4D_0r=4D_0s_0$.
\end{proof}

\begin{lemma}[Primary scale limits]
    \label{thm:area_law_primary_scale_limits}
    \lean{TNLean.PEPS.AreaLaw.Geometry.eventually_primaryScaleLimits,
        TNLean.PEPS.AreaLaw.Geometry.le_primaryScannerScale}
    \leanok
    For any $C_P$ and $\varepsilon>0$, all sufficiently large $k$ (uniformly
    over primaries) satisfy
    \begin{align}
        \bigl\lfloor n_k^{1-\ell}\bigr\rfloor & \le\varepsilon r,   &
        n_k^\beta                             & \le\varepsilon t_k, &
        n_k^\beta                             & \le\varepsilon r,
        \label{eq:area_law_primary_scale_limits}
    \end{align}
    and if $C_P>0$, then for every $n_*$, $n_k\ge n_*$ for all sufficiently
    large $k$.
\end{lemma}
\begin{proof}\leanok
    \uses{def:area_law_primary_birth_template,thm:area_law_exponent_gaps}
    With $e=1+2\delta_0$ and $C=\max\{C_P,0\}+1$,~(\ref{eq:area_law_primary_scanner_scale})
    gives $n_k\le Cr^e$. Since $e(1-\ell)<1$, $e\beta<\zeta$
    (Lemma~\ref{thm:area_law_exponent_gaps}), and $r^\zeta/2\le t_k\le r$,
    $Cr^{-(b-a)}\to0$ ($a<b$) yields~(\ref{eq:area_law_primary_scale_limits});
    for $C_P>0$, $n_k\ge C_P r\to\infty$.
\end{proof}
